\section{Grasp sources and simulator mapping}
\subsection{Validation framework}
Each source record identifies its input geometry, physical scale, coordinate
frames, hand assets, checkpoint, configuration and output stage. We first
reconstruct the native proposal and independently check transformation,
joint-order and forward-kinematic consistency. We then compare native and
translated configurations in a shared object frame, reporting per-finger
landmark errors, object proximity, penetration, lost and added near-surface
regions, and reliable surface-direction comparisons. Visual and collision
surfaces are distinct references, with unmatched regions and sampling uncertainty
retained in every report.

Numerical decoding, geometric preservation and execution adequacy are separate
claims. Agreement of joints or fingertips does not establish equivalent
collision surfaces, contact forces or task outcomes. Poor native proposals
remain in the audit; successful lifting is not a mapping acceptance criterion.
Thresholds and sampling assumptions are declared explicitly, and inspected
development cases are distinguished from later independent validation.
Our common executor differs from source-specific policies and quality tests.

For the current HUG development audit, a frozen strict encoded-surface claim
is refuted by a positive same-part lost/added proximity-area lower bound or a
reliable object-facing sign reversal, under declared numerical assumptions.
An absent sampled reversal cannot certify continuous normal agreement. That
claim defines strict preservation, not a universal tolerance for all useful
retargeters. An approximate native-fidelity claim requires a separately
justified distortion budget and independent evidence. By contrast, checking
that a declared retargeter and simulator realize the specified Shadow pose
does not require native outcome equivalence or a universal morphology
allowance. Studying that executed pipeline still requires distortion reporting,
mapper/geometry confounding controls and independent prediction evidence.
None of these narrower claims retroactively turns the recorded HUG strict
failures into passes.

Provider audits and their evidence-backed manuscript subsections can proceed
in parallel. Shared retargeters, assets, measurement code and execution settings
are versioned dependencies: a change requires reviewing affected audits.
The current phase uses static reconstruction of saved inputs and manufactured
checks. It generates no new evaluation episodes, physical outcome labels or
training data. Actual-native validation remains pending where authentic source
outputs or assets have not been established.
